\documentclass[border=8pt, multi, tikz]{standalone}
\usepackage{import}
\usepackage{tikz}
\usetikzlibrary{positioning, arrows.meta, shapes.geometric, calc, fit}

\definecolor{inputcolor}{RGB}{220,220,255}
\definecolor{conv1color}{RGB}{255,230,150}
\definecolor{conv2color}{RGB}{255,200,120}
\definecolor{conv3color}{RGB}{255,170,90}
\definecolor{latentcolor}{RGB}{255,100,100}
\definecolor{densecolor}{RGB}{150,255,150}
\definecolor{outputcolor}{RGB}{220,220,255}

\tikzset{
    block/.style={
        rectangle,
        draw,
        thick,
        minimum width=2.5cm,
        minimum height=1.5cm,
        align=center,
        font=\Large\bfseries,
        drop shadow
    },
    arrow/.style={
        -Stealth,
        very thick,
        draw=black!80
    },
    label/.style={
        font=\large,
        align=center
    }
}

\begin{document}
\begin{tikzpicture}[node distance=2.5cm, auto]

    % Title
    \node[font=\Huge\bfseries] at (0, 12) {Bowhead Whale Autoencoder Architecture};
    \node[font=\Large] at (0, 11) {Hybrid 5×5/3×3 Convolutional Autoencoder};
    
    % ========== ENCODER ==========
    \node[font=\LARGE\bfseries, text=blue!70!black] at (-8, 9.5) {ENCODER};
    
    % Input
    \node[block, fill=inputcolor, minimum height=2cm] (input) at (-8, 7.5) {
        \textbf{Input}\\[0.3cm]
        1×121×104\\[0.1cm]
        \small Spectrogram
    };
    
    % Conv1 + Pool1
    \node[block, fill=conv1color, right=of input] (conv1) {
        \textbf{Conv1}\\[0.2cm]
        5×5 kernel\\
        32 filters\\[0.1cm]
        \small +BN+ReLU+Pool
    };
    \node[label, below=0.1cm of conv1] {\small Output: 32×60×52};
    
    % Conv2 + Pool2
    \node[block, fill=conv2color, right=of conv1] (conv2) {
        \textbf{Conv2}\\[0.2cm]
        3×3 kernel\\
        64 filters\\[0.1cm]
        \small +BN+ReLU+Pool
    };
    \node[label, below=0.1cm of conv2] {\small Output: 64×30×26};
    
    % Conv3 + Pool3
    \node[block, fill=conv3color, right=of conv2] (conv3) {
        \textbf{Conv3}\\[0.2cm]
        3×3 kernel\\
        128 filters\\[0.1cm]
        \small +BN+ReLU+Pool
    };
    \node[label, below=0.1cm of conv3] {\small Output: 128×15×13};
    
    % ========== BOTTLENECK ==========
    \node[font=\LARGE\bfseries, text=red!70!black] at (4, 4) {BOTTLENECK};
    
    % Flatten
    \node[block, fill=densecolor, below=of conv3, minimum width=2cm] (flatten) {
        \textbf{Flatten}\\[0.2cm]
        26,624
    };
    
    % Dense 1
    \node[block, fill=densecolor, below=of flatten, minimum width=2cm] (dense1) {
        \textbf{Dense}\\[0.2cm]
        64\\[0.1cm]
        \small +ReLU
    };
    
    % Latent Space
    \node[block, fill=latentcolor, below=of dense1, minimum width=3cm, minimum height=2cm] (latent) {
        \textbf{LATENT}\\
        \textbf{SPACE}\\[0.3cm]
        \Huge 32-dim
    };
    
    % ========== DECODER ==========
    \node[font=\LARGE\bfseries, text=blue!70!black] at (8, -5) {DECODER};
    
    % Dense 2
    \node[block, fill=densecolor, below=of latent, minimum width=2cm] (dense2) {
        \textbf{Dense}\\[0.2cm]
        64\\[0.1cm]
        \small +ReLU
    };
    
    % Dense 3
    \node[block, fill=densecolor, below=of dense2, minimum width=2cm] (dense3) {
        \textbf{Dense}\\[0.2cm]
        26,624\\[0.1cm]
        \small +ReLU
    };
    
    % Reshape
    \node[block, fill=densecolor, below=of dense3, minimum width=2.5cm] (reshape) {
        \textbf{Reshape}\\[0.2cm]
        128×15×13
    };
    
    % Deconv1
    \node[block, fill=conv3color, left=of reshape] (deconv1) {
        \textbf{TransConv1}\\[0.2cm]
        2×2 kernel\\
        64 filters\\[0.1cm]
        \small +BN+ReLU
    };
    \node[label, below=0.1cm of deconv1] {\small Output: 64×30×26};
    
    % Deconv2
    \node[block, fill=conv2color, left=of deconv1] (deconv2) {
        \textbf{TransConv2}\\[0.2cm]
        2×2 kernel\\
        32 filters\\[0.1cm]
        \small +BN+ReLU
    };
    \node[label, below=0.1cm of deconv2] {\small Output: 32×60×52};
    
    % Deconv3 (Output)
    \node[block, fill=conv1color, left=of deconv2] (deconv3) {
        \textbf{TransConv3}\\[0.2cm]
        2×2 kernel\\
        1 filter
    };
    \node[label, below=0.1cm of deconv3] {\small Output: 1×121×104};
    
    % Output
    \node[block, fill=outputcolor, left=of deconv3, minimum height=2cm] (output) {
        \textbf{Output}\\[0.3cm]
        1×121×104\\[0.1cm]
        \small Reconstruction
    };
    
    % ========== ARROWS ==========
    \draw[arrow] (input) -- (conv1);
    \draw[arrow] (conv1) -- (conv2);
    \draw[arrow] (conv2) -- (conv3);
    \draw[arrow] (conv3) -- (flatten);
    \draw[arrow] (flatten) -- (dense1);
    \draw[arrow] (dense1) -- (latent);
    \draw[arrow] (latent) -- (dense2);
    \draw[arrow] (dense2) -- (dense3);
    \draw[arrow] (dense3) -- (reshape);
    \draw[arrow] (reshape) -- (deconv1);
    \draw[arrow] (deconv1) -- (deconv2);
    \draw[arrow] (deconv2) -- (deconv3);
    \draw[arrow] (deconv3) -- (output);
    
    % ========== ANNOTATIONS ==========
    % Architecture type
    \node[draw, thick, rounded corners, fill=yellow!20, font=\large, align=center] at (0, -14) {
        \textbf{Hybrid Architecture:}\\
        First layer uses 5×5 kernel to capture broad N/U-shaped whale call patterns\\
        Subsequent layers use 3×3 kernels for efficient detail refinement
    };
    
    % Model stats
    \node[draw, thick, rounded corners, fill=green!10, font=\large, align=left] at (-8, -16.5) {
        \textbf{Model Statistics:}\\
        • Parameters: 3,359,105\\
        • Latent Dimension: 32\\
        • Compression Ratio: 378:1
    };
    
    % Training config
    \node[draw, thick, rounded corners, fill=blue!10, font=\large, align=left] at (8, -16.5) {
        \textbf{Training Configuration:}\\
        • Optimizer: Adam (lr=1e-3)\\
        • Loss: MSE Reconstruction\\
        • Batch Normalization: Enabled
    };

\end{tikzpicture}
\end{document}
